% ============================================================
% ÉVALUATION — STS BIOALC 2e ANNÉE
% Partie 1 — Séance du lundi (50 min)
% Durée : 30 minutes — Barème : /20
% ============================================================

\documentclass[11pt,a4paper]{article}

\usepackage[utf8]{inputenc}
\usepackage[T1]{fontenc}
\usepackage[french]{babel}
\usepackage{amsmath,amssymb}
\usepackage{geometry}
\usepackage{enumitem}
\usepackage{booktabs}
\usepackage{array}
\usepackage{tabularx}
\usepackage{xcolor}
\usepackage{tcolorbox}
\usepackage{fancyhdr}
\usepackage{multicol}
\usepackage{tikz}
\usepackage{pgfplots}
\pgfplotsset{compat=1.18}
\usepackage{lastpage}

\tcbuselibrary{skins,breakable}

\geometry{margin=2cm, headheight=15pt}

% ------------------------------------------------------------
% Couleurs
% ------------------------------------------------------------
\definecolor{primary}{RGB}{31,78,121}
\definecolor{soft}{RGB}{230,239,247}
\definecolor{amberbg}{RGB}{255,244,224}
\definecolor{amberborder}{RGB}{240,192,112}
\definecolor{greenbg}{RGB}{233,246,234}
\definecolor{greenborder}{RGB}{182,224,184}
\definecolor{redbg}{RGB}{253,236,236}
\definecolor{redborder}{RGB}{242,179,171}
\definecolor{grayline}{RGB}{120,130,140}

% ------------------------------------------------------------
% En-tête et pied de page
% ------------------------------------------------------------
\pagestyle{fancy}
\fancyhf{}
\fancyhead[L]{\textcolor{primary}{\textbf{BTS BioALC — 2\textsuperscript{ème} année}}}
\fancyhead[R]{\textcolor{primary}{Évaluation — Partie 1}}
\fancyfoot[L]{\textcolor{grayline}{\small Nom : \dotfill}}
\fancyfoot[C]{\textcolor{grayline}{\small \thepage/\pageref{LastPage}}}
\fancyfoot[R]{\textcolor{grayline}{\small Durée : 30 min}}
\renewcommand{\headrulewidth}{0.4pt}
\renewcommand{\footrulewidth}{0.4pt}

% ------------------------------------------------------------
% Boîtes
% ------------------------------------------------------------
\newtcolorbox{consigne}[1][]{
  colback=soft, colframe=primary, boxrule=0.8pt, arc=4pt,
  fonttitle=\bfseries, coltitle=white, colbacktitle=primary,
  left=6pt, right=6pt, top=4pt, bottom=4pt,
  breakable, #1
}

\newtcolorbox{contexte}{
  colback=amberbg, colframe=amberborder, boxrule=0.8pt, arc=4pt,
  title={\textbf{Contexte professionnel}}, fonttitle=\bfseries,
  coltitle=black, colbacktitle=amberborder,
  left=6pt, right=6pt, top=4pt, bottom=4pt,
  breakable
}

\newtcolorbox{questionbox}{
  colback=white, colframe=primary, boxrule=0.8pt, arc=4pt,
  fonttitle=\bfseries, coltitle=white, colbacktitle=primary,
  left=6pt, right=6pt, top=4pt, bottom=4pt,
  breakable
}

% ------------------------------------------------------------
% Commande lignes d'écriture
% ------------------------------------------------------------
\newcommand{\lignes}[1]{%
  \vspace{0.15cm}%
  \foreach \x in {1,...,#1}{%
    \noindent\rule{\linewidth}{0.4pt}\par\vspace{0.55cm}%
  }%
}

\title{}
\author{}
\date{}

% ============================================================
\begin{document}
% ============================================================

\begin{center}
  {\Large \textbf{Évaluation — Statistique descriptive}}\\[0.2cm]
  {\large \textbf{Partie 1 — Séance du lundi (50 min)}}\\[0.2cm]
  {\large Durée : 30 minutes \hfill Barème : \textbf{/20}}
\end{center}

\vspace{0.3cm}

\textbf{Nom :} \dotfill \quad \textbf{Prénom :} \dotfill \quad \textbf{Date :} \dotfill

\vspace{0.4cm}

\begin{consigne}
\textbf{Consignes :}
\begin{itemize}[leftmargin=*]
  \item Durée : 30 minutes.
  \item Calculatrice autorisée. Documents interdits.
  \item Répondre directement sur le sujet.
  \item La qualité de la rédaction, les justifications et les unités seront prises en compte.
  \item Cette évaluation porte uniquement sur la \textbf{Partie 1 — Séance du lundi}.
  \item \textbf{Répartition indicative du temps :} Exercice 1 : 5 min — Exercice 2 : 15 min — Exercice 3 : 10 min.
\end{itemize}
\end{consigne}

\vspace{0.3cm}

\begin{contexte}
Vous êtes technicien(ne) dans un laboratoire de microbiologie. Dans le cadre d’un contrôle de croissance bactérienne, le responsable vous demande d’analyser le \textbf{diamètre de colonies} (en mm) obtenues sur \textbf{20 boîtes de Petri} après 24~h d’incubation à 37~°C. Votre mission : produire un tableau de synthèse et un graphique permettant de caractériser la dispersion des tailles de colonies observées.
\end{contexte}

\vspace{0.4cm}

% ============================================================
\section*{Exercice 1 — Restitution de connaissances \hfill \normalfont\normalsize(4 points)}
% ============================================================

\begin{questionbox}
\begin{enumerate}[leftmargin=*]
  \item Qu'appelle-t-on \textbf{population} en statistique ? Et \textbf{échantillon} ? \hfill \textbf{(2 pts)}
  \lignes{4}
  \item Quelle est la différence entre une variable \textbf{qualitative} et une variable \textbf{quantitative} ? Donner un exemple de chaque. \hfill \textbf{(2 pts)}
  \lignes{4}
\end{enumerate}
\end{questionbox}

% ============================================================
\newpage
% ============================================================

\section*{Exercice 2 — Tableau et représentation graphique \hfill \normalfont\normalsize(10 points)}

\begin{questionbox}
Voici les diamètres de colonies (en mm) mesurés sur 20 boîtes de Petri :

\begin{center}
\begin{tabular}{|c|c|c|c|c|c|c|c|c|c|}
\hline
5,2 & 6,8 & 4,5 & 7,1 & 5,9 & 6,3 & 4,8 & 7,5 & 5,4 & 6,1 \\
\hline
4,9 & 6,7 & 5,7 & 4,3 & 7,2 & 6,0 & 5,5 & 6,9 & 4,7 & 6,4 \\
\hline
\end{tabular}
\end{center}

\vspace{0.4cm}

\textbf{Question 1.} Regroupez ces données en \textbf{classes de 1 mm} (de 4 à 8 mm) et complétez le tableau ci-dessous. \hfill \textbf{(5 pts)}

\vspace{0.3cm}
{\renewcommand{\arraystretch}{1.8}
\begin{center}
\begin{tabular}{|c|c|c|c|}
\hline
\textbf{Classe (mm)} & \textbf{Effectif $n_i$} & \textbf{Fréquence $f_i$} & \textbf{Fréquence cumulée croissante} \\
\hline
[4 ; 5[ & & & \\
\hline
[5 ; 6[ & & & \\
\hline
[6 ; 7[ & & & \\
\hline
[7 ; 8] & & & \\
\hline
\textbf{Total} & & & \\
\hline
\end{tabular}
\end{center}
}

\vspace{0.4cm}

\textbf{Question 2.} Quel type de graphique est adapté à cette variable ? Justifiez votre réponse. \hfill \textbf{(2 pts)}

\lignes{6}
\end{questionbox}

% ============================================================
\newpage
% ============================================================

\begin{questionbox}
\textbf{Question 3.} Construisez ce graphique dans le cadre ci-dessous. Pensez à indiquer le titre, les axes et les unités. \hfill \textbf{(3 pts)}

\vspace{0.3cm}
\begin{center}
\begin{tikzpicture}
% Grille en orientation portrait (9 cm de large × 16 cm de haut)
\draw[gray!50, step=0.5cm] (0,0) grid (9,16);
% Axe horizontal (bas) : effectifs
\draw[thick, ->] (0,0) -- (9.3,0);
% Axe vertical (gauche) : classes de diamètre
\draw[thick, ->] (0,0) -- (0,16.3);
% Étiquette de l'axe horizontal (effectifs)
\node at (4.5,-0.8) {\small Effectif};
% Étiquette de l'axe vertical (diamètre des colonies)
\node[rotate=90] at (-1.2,8) {\small Diamètre des colonies (mm)};
\end{tikzpicture}
\end{center}
\end{questionbox}

% ============================================================
\newpage
% ============================================================

\section*{Exercice 3 — Synthèse et analyse critique \hfill \normalfont\normalsize(6 points)}

\begin{questionbox}
\begin{enumerate}[leftmargin=*]
  \item Pour une variable \textbf{qualitative}, quel type de graphique utiliser ? \hfill \textbf{(1,5 pt)}
  \lignes{2}

  \item À partir du tableau complété à l’exercice 2 : \hfill \textbf{(2,5 pts)}
        \begin{enumerate}[label=(\alph*), leftmargin=*]
          \item Quelle est la valeur de la \textbf{fréquence cumulée croissante} correspondant à la classe [5 ; 6[ ? \hfill \textbf{(1 pt)}
          \lignes{2}
          \item Interprétez ce résultat en une phrase, dans le contexte de l’étude. \hfill \textbf{(1,5 pt)}
          \lignes{4}
        \end{enumerate}

  \item \textbf{Analyse critique.} On vous présente le graphique ci-dessous, représentant les effectifs par classe de diamètre de colonies. \hfill \textbf{(2 pts)}

  \vspace{0.3cm}
  \begin{center}
  \begin{tikzpicture}[scale=1.0]
    % Axes
    \draw[thick, ->] (0,0) -- (5.5,0);
    \draw[thick, ->] (0,0) -- (0,9.5);
    % Barres (effectifs : 5, 5, 7, 3)
    \draw[fill=primary!20] (0.7,0) rectangle (1.3,5);
    \draw[fill=primary!20] (1.7,0) rectangle (2.3,5);
    \draw[fill=primary!20] (2.7,0) rectangle (3.3,7);
    \draw[fill=primary!20] (3.7,0) rectangle (4.3,3);
    % Étiquettes axe X (classes sans unité)
    \node at (1, -0.6) {\small [4;5[};
    \node at (2, -0.6) {\small [5;6[};
    \node at (3, -0.6) {\small [6;7[};
    \node at (4, -0.6) {\small [7;8]};
    % Graduations axe Y (échelle non régulière)
    \foreach \y/\label in {0/0, 1/1, 2/2, 3/4, 4/8} {
      \draw (0,\y) -- (-0.1,\y) node[left] {\small \label};
    }
    % Pas de titre, pas d'unités sur les axes
  \end{tikzpicture}
  \end{center}

  \vspace{0.3cm}
  \textbf{Quelles sont les erreurs liées à ce graphique ?} \\
  Citez-en deux et proposez une correction pour chacune directement sur le graphique ci-dessous.
  \lignes{6}
   \vspace{0.3cm}
  \begin{center}
  	\begin{tikzpicture}[scale=1.0]
  		% Axes
  		\draw[thick, ->] (0,0) -- (5.5,0);
  		\draw[thick, ->] (0,0) -- (0,9.5);
  		% Barres (effectifs : 5, 5, 7, 3)
  		\draw[fill=primary!20] (0.7,0) rectangle (1.3,5);
  		\draw[fill=primary!20] (1.7,0) rectangle (2.3,5);
  		\draw[fill=primary!20] (2.7,0) rectangle (3.3,7);
  		\draw[fill=primary!20] (3.7,0) rectangle (4.3,3);
  		% Étiquettes axe X (classes sans unité)
  		\node at (1, -0.6) {\small [4;5[};
  		\node at (2, -0.6) {\small [5;6[};
  		\node at (3, -0.6) {\small [6;7[};
  		\node at (4, -0.6) {\small [7;8]};
  		% Graduations axe Y (échelle non régulière)
  		\foreach \y/\label in {0/0, 1/1, 2/2, 3/4, 4/8} {
  			\draw (0,\y) -- (-0.1,\y) node[left] {\small \label};
  		}
  		% Pas de titre, pas d'unités sur les axes
  	\end{tikzpicture}
  \end{center}
\end{enumerate}
\end{questionbox}

\vfill

% ============================================================
% Récapitulatif du barème
% ============================================================


\vspace{0.3cm}

\begin{center}
\textit{Évaluation — BTS BioALC — 2\textsuperscript{ème} année — Partie 1 — Semaine du 05/10/2026}
\end{center}

\end{document}